\section{Grouped Contextual Mobile Panels}
\label{sec:pilot-mobile-contextual-panels}

The unified Pilot phone interface now presents one working context at a time.  This
completes mobile task \texttt{MOB-0212} without removing any Personal, Workspace or
Boardroom capability.  The change replaces the earlier cumulative controller stream,
in which tools, inbox items, action cards and connection settings could remain visible
in one long page, with a small routing shell whose visible body is determined by the
active product mode, bottom-navigation destination and selected tool.

\subsection{Interaction model}

The user first chooses Personal, Workspace or Boardroom.  Within Pilot, a grouped
selector exposes the complete signed shortcut manifest.  Personal tools are grouped
as Everyday, Home and experiences, and Private life; Workspace tools are grouped as
Assignment and Resources; Boardroom tools are grouped as Business and People and
resources.  Four high-frequency Personal controls---TV, Tasks, Calendar and
Devices---remain directly reachable without opening the selector.

Selecting a tool replaces the overview with that tool's panel.  A visible ``Close
tool'' control returns to the relevant mode overview.  Selecting Inbox, Actions,
Spaces or Me clears the former tool selection before rendering the new destination.
Consequently private Memory or Guardian controls cannot remain underneath an Inbox or
connection screen, and generic cards are shown only in the Pilot overview or Actions
context.

The persistent Pilot composer and primary navigation remain available independently
of the selected panel.  This preserves the product's conversational centre while
preventing advanced controls from becoming a permanently expanded settings page.

\subsection{Deterministic rendering boundary}

Let $m$ be the active product mode, $n$ the active primary-navigation destination and
$t$ the optional selected tool.  The visible capability panel is
\[
  P(m,n,t)=
  \left\{\begin{array}{ll}
    \mathrm{Tool}(m,t), & n=\texttt{Pilot}\ \land\ t\neq\emptyset,\\
    \mathrm{Overview}(m), & n=\texttt{Pilot}\ \land\ t=\emptyset,\\
    \mathrm{Destination}(n), & n\in\{\texttt{Inbox},\texttt{Actions},
      \texttt{Spaces},\texttt{Me}\}.
  \end{array}\right.
\]
Changing mode resets $n$ to \texttt{Pilot} and clears $t$.  Leaving Pilot clears $t$.
Each sensitive Personal panel additionally requires a valid phone connection and the
existing mother-brain capability check; hiding or revealing a panel never grants
authority.

\subsection{Security and privacy properties}

The restructuring is a presentation change over existing governed endpoints.  It
does not move credentials, memory, contact data, messages or business records into
the browser.  The mother brain continues to validate identity, lease, scope, exact
approval and execution authority.  The client does not infer capability from a
visible option, and a stale selected tool cannot survive a transition into another
primary context.

The selector is populated from the signed surface manifest rather than a household-
specific hard-coded menu.  The same shell therefore remains suitable for an owner,
employee, freelancer, child or Boardroom member while showing only the tools supplied
for that person's active space.

\subsection{Verification and remaining field gate}

Automated checks confirm grouped selector semantics, manifest-defined option groups,
context isolation, Memory and Guardian isolation, overview-card boundaries and direct
access to TV, Tasks, Calendar and Devices.  The optimized mobile production build
type-checks and renders the revised shell successfully.

Formal Stage~2 closure still requires \texttt{MOB-0216}: representative small and
large iPhone and Android viewport testing with real touch, dynamic type and assistive
technology.  That device-field evidence is deliberately left open and must not be
inferred from browser compilation or source-level tests.
